\documentclass{csbeamer}

\usepackage{hyperref}
\usepackage{listings}
\usepackage{xcolor}
\usepackage{graphicx}
\usepackage{tikz}
\usetikzlibrary{positioning,shapes,arrows,calc}

% Python code styling
\lstset{
    language=Python,
    basicstyle=\ttfamily\scriptsize,
    keywordstyle=\color{blue}\bfseries,
    stringstyle=\color{red},
    commentstyle=\color{green!50!black}\itshape,
    showstringspaces=false,
    breaklines=true,
    frame=single,
    numbers=left,
    numberstyle=\tiny\color{gray}
}

% Course information
\title{Tabular Data and CSV Files}
\course{CSCI 128: Coding for Problem Solving}
\courseshort{CSCI 128}
\term{Winter 2026}
\author{Dr. Jean-Alexis Delamer}

\begin{document}

\begin{frame}
\titlepage
\end{frame}

% -----------------------------------------------
\begin{frame}{Today's Topics}
    \begin{center}
        \Large From pixels to spreadsheets!
    \end{center}

    \vspace{1em}

    \textbf{Learning Objectives:}
    \begin{itemize}
        \item<1-> Understand \textcolor{stfxblue}{tabular data} as a programming concept
        \item<2-> Read and write \textcolor{marigold}{CSV files} in Python
        \item<3-> Access rows and columns by \textcolor{stfxblue}{index}
        \item<4-> Compute basic statistics: \textcolor{marigold}{min}, \textcolor{marigold}{max}, \textcolor{marigold}{average}
        \item<5-> Filter rows based on \textcolor{stfxblue}{conditions}
        \item<6-> Process real datasets with \textcolor{marigold}{loops}
        \item<7-> Build reusable \textcolor{stfxblue}{data-processing functions}
    \end{itemize}

    \vspace{1em}

    \begin{center}
        \onslide<8->{\textit{Tables are just \textcolor{marigold}{\textbf{grids of data}}, and grids are something we already know!}}
    \end{center}
\end{frame}

% -----------------------------------------------
\begin{frame}{From Images to Tables}
    \begin{center}
        \textbf{You already know 2D grids, this is the same idea!}
    \end{center}

    \vspace{0.5em}

    \begin{center}
        \begin{tabular}{l|l}
            \textbf{\textcolor{stfxblue}{Images}} & \textbf{\textcolor{marigold}{Tables}} \\
            \hline
            2D grid of \textcolor{stfxblue}{pixels}    & 2D grid of \textcolor{marigold}{values} \\
            Pixel = (R, G, B)                           & Cell = a number or text \\
            Width $\times$ Height                       & Columns $\times$ Rows \\
            Modify pixels                               & Modify cells \\
            Filter by color                             & Filter by value \\
        \end{tabular}
    \end{center}

    \vspace{1em}

    \begin{center}
        \onslide<2->{\Large \textcolor{marigold}{\textbf{Both are just grids of data!}}}
    \end{center}

    \vspace{0.5em}

    \onslide<3->{
    \begin{center}
        \textit{The skills you built with images transfer directly to tables.}
    \end{center}
    }
\end{frame}

% -----------------------------------------------
\begin{frame}[fragile]{What is a CSV File?}
    \textbf{CSV = Comma-Separated Values}

    \vspace{0.5em}

    \begin{itemize}
        \item A plain text file, you can open it in Notepad
        \item Each \textcolor{stfxblue}{line} is one row of data
        \item \textcolor{marigold}{Commas} separate the columns within each row
        \item The \textcolor{stfxblue}{first row} is usually a \textbf{header} (column names)
        \item Universally supported: Excel, Python, R, databases \ldots
    \end{itemize}

    \vspace{0.5em}

    \textbf{Example: \texttt{students.csv}}
    \begin{lstlisting}[language={},frame=single,numbers=none]
Name,Grade,Score
Alice,A,92
Bob,B,78
Carol,A,95
David,C,64
Eve,B,81
    \end{lstlisting}
\end{frame}

% -----------------------------------------------
\begin{frame}[fragile]{Reading a CSV: Step 1 - Open the File}
    \textbf{First, just open the file and look at the raw lines:}

    \begin{lstlisting}
with open("students.csv", "r") as f:
    lines = f.readlines()

print(lines)
    \end{lstlisting}

    \vspace{0.8em}
    \pause

    \textbf{Output - each line is one string, including the newline \texttt{\textbackslash n}:}

    \begin{lstlisting}
['Name,Grade,Score\n', 'Alice,A,92\n', 'Bob,B,78\n', ...]
    \end{lstlisting}

    \vspace{0.8em}
    \pause

    \textbf{Key ideas:}
    \begin{itemize}
        \item<3-> \texttt{open("file", "r")} opens a file for \textcolor{stfxblue}{reading}
        \item<4-> \texttt{f.readlines()} returns a \textcolor{marigold}{\textbf{list of strings}}, one per line
        \item<5-> Every line still has a \texttt{\textbackslash n} (newline) at the end - we need to clean that up
    \end{itemize}
\end{frame}

% -----------------------------------------------
\begin{frame}[fragile]{Reading a CSV: Step 2 - Clean Up Each Line}
    \textbf{The \texttt{.strip()} method removes whitespace and newlines:}

    \begin{lstlisting}
line = "Alice,A,92\n"

clean = line.strip()
print(clean)   # Alice,A,92
    \end{lstlisting}

    \vspace{0.5em}
    \pause

    \textbf{The \texttt{.split(",")} method breaks a string on commas:}

    \begin{lstlisting}
clean = "Alice,A,92"

parts = clean.split(",")
print(parts)   # ['Alice', 'A', '92']
    \end{lstlisting}

    \vspace{0.5em}
    \pause

    \textbf{Chain them together:}

    \begin{lstlisting}
row = "Alice,A,92\n".strip().split(",")
print(row)     # ['Alice', 'A', '92']
    \end{lstlisting}
\end{frame}

% -----------------------------------------------
\begin{frame}[fragile]{Reading a CSV: Step 3 - Read the Header}
    \textbf{The \textcolor{stfxblue}{first line} is special - it contains the column names:}

    \begin{lstlisting}
with open("students.csv", "r") as f:
    lines = f.readlines()

# lines[0] is the first line (the header)
header = lines[0].strip().split(",")
print("Columns:", header)
# Columns: ['Name', 'Grade', 'Score']
    \end{lstlisting}

    \pause

    \textbf{The remaining lines (\texttt{lines[1:]}) are the data rows:}

    \begin{lstlisting}
print("Number of data rows:", len(lines[1:])) # Number of data rows: 5

print("First data row (raw):", lines[1]) # First data row (raw): Alice,A,92
    \end{lstlisting}

\end{frame}

% -----------------------------------------------
\begin{frame}[fragile]{Reading a CSV: Step 4 - Parse All Data Rows}
    \textbf{Loop through \texttt{lines[1:]} and split each one:}

    \begin{lstlisting}
with open("students.csv", "r") as f:
    lines = f.readlines()

data = []
for line in lines[1:]:       # skip the header (index 0)
    line = line.strip()      # remove the newline character
    if line == "":           # skip any accidental blank lines
        continue
    row = line.split(",")    # split on commas -> a list
    data.append(row)         # add the list to data

print(data)
# [['Alice', 'A', '92'], ['Bob', 'B', '78'], ...]
    \end{lstlisting}

    \pause

    \textbf{\texttt{data} is now a \textcolor{marigold}{list of lists} - one inner list per student.}
\end{frame}

% -----------------------------------------------
\begin{frame}[fragile]{Reading a CSV: Step 5 - Access Individual Cells}
    \textbf{Each row is a list: \texttt{row[0]} = Name, \texttt{row[1]} = Grade, \texttt{row[2]} = Score}

    \begin{lstlisting}
for row in data:
    name  = row[0]
    grade = row[1]
    score = row[2]
    print(name, "got a", grade, "with score", score)
    \end{lstlisting}

    \pause

    \textbf{Output:}

    \begin{lstlisting}
Alice got a A with score 92
Bob got a B with score 78
Carol got a A with score 95
David got a C with score 64
Eve got a B with score 81
    \end{lstlisting}

    \pause

    \begin{itemize}
        \item<3-> Scores are \textcolor{stfxblue}{\textbf{strings}}, not numbers - \texttt{"92"}, not \texttt{92}
        \item<4-> To do math: \texttt{int(row[2])} or \texttt{float(row[2])}
    \end{itemize}
\end{frame}

% -----------------------------------------------
\begin{frame}[fragile]{Reading a CSV File: Using the \texttt{csv} Module}
    \textbf{Python's standard library has a \textcolor{stfxblue}{\texttt{csv}} module - no install needed!}

    \begin{lstlisting}
import csv

# csv.reader handles commas, quotes, and newlines for us
with open("students.csv", "r", newline="") as f:
    reader = csv.reader(f)

    header = next(reader)       # read the first row as header
    print("Columns:", header)   # Columns: ['Name', 'Grade', 'Score']

    data = []
    for row in reader:          # each row is already a list
        data.append(row)
    \end{lstlisting}

\end{frame}

% -----------------------------------------------
\begin{frame}[fragile]{The \texttt{csv} Module: Key Takeaways}
    \textbf{Why use \texttt{csv.reader} instead of doing it manually?}

    \vspace{0.8em}

    \begin{itemize}
        \item<1-> \textcolor{stfxblue}{\textbf{Simpler code}} - no \texttt{.strip()}, no \texttt{.split(",")}, no blank-line checks
        \item<2-> \textcolor{marigold}{\textbf{Handles edge cases}} - e.g.\ values that contain commas inside quotes:\\
              \texttt{"Smith, John",A,92} is parsed correctly
        \item<3-> Each row comes back as a \textcolor{stfxblue}{\textbf{list of strings}} - just like the manual way
        \item<4-> \texttt{next(reader)} reads \textcolor{marigold}{\textbf{one row}} and advances the cursor - useful for skipping the header
    \end{itemize}

    \vspace{0.8em}

    \pause

    \begin{block}{Remember: all values are strings}
        To do arithmetic, convert first:\quad
        \texttt{int(row[2])} gives \texttt{92}\quad
        \texttt{float(row[2])} gives \texttt{92.0}
    \end{block}
\end{frame}

% -----------------------------------------------
\begin{frame}{Understanding Rows and Columns}
    \textbf{Think of your table like a 2D grid - just like an image:}

    \vspace{0.5em}

    \begin{center}
        \begin{tikzpicture}[scale=0.85]
            % Column header labels
            \node[font=\bfseries\small, text=stfxblue] at (2.0, 3.6) {col 0 (\texttt{Name})};
            \node[font=\bfseries\small, text=stfxblue] at (4.5, 3.6) {col 1 (\texttt{Grade})};
            \node[font=\bfseries\small, text=stfxblue] at (7.0, 3.6) {col 2 (\texttt{Score})};

            % Row 0
            \draw[fill=blue!10] (0.7,2.8) rectangle (3.3,3.2);
            \draw[fill=blue!10] (3.3,2.8) rectangle (5.7,3.2);
            \draw[fill=blue!10] (5.7,2.8) rectangle (8.3,3.2);
            \node at (2.0,3.0) {\small Alice};
            \node at (4.5,3.0) {\small A};
            \node at (7.0,3.0) {\small 92};
            \node[font=\bfseries\small] at (0.1,3.0) {row 0};

            % Row 1
            \draw[fill=green!8] (0.7,2.4) rectangle (3.3,2.8);
            \draw[fill=green!8] (3.3,2.4) rectangle (5.7,2.8);
            \draw[fill=green!8] (5.7,2.4) rectangle (8.3,2.8);
            \node at (2.0,2.6) {\small Bob};
            \node at (4.5,2.6) {\small B};
            \node at (7.0,2.6) {\small 78};
            \node[font=\bfseries\small] at (0.1,2.6) {row 1};

            % Row 2
            \draw[fill=blue!10] (0.7,2.0) rectangle (3.3,2.4);
            \draw[fill=blue!10] (3.3,2.0) rectangle (5.7,2.4);
            \draw[fill=blue!10] (5.7,2.0) rectangle (8.3,2.4);
            \node at (2.0,2.2) {\small Carol};
            \node at (4.5,2.2) {\small A};
            \node at (7.0,2.2) {\small 95};
            \node[font=\bfseries\small] at (0.1,2.2) {row 2};

            % Outer border
            \draw[thick] (0.7,2.0) rectangle (8.3,3.2);
        \end{tikzpicture}
    \end{center}

    \vspace{0.5em}

    \textbf{Accessing values:}
    \begin{itemize}
        \item<2-> \texttt{data[0]}        \hfill $\rightarrow$ \texttt{['Alice', 'A', '92']}  (whole row)
        \item<3-> \texttt{data[0][0]}     \hfill $\rightarrow$ \texttt{'Alice'}               (row 0, col 0)
        \item<4-> \texttt{data[1][2]}     \hfill $\rightarrow$ \texttt{'78'}                  (row 1, col 2)
        \item<5-> \texttt{data[2][2]}     \hfill $\rightarrow$ \texttt{'95'}                  (row 2, col 2)
    \end{itemize}
\end{frame}

% -----------------------------------------------
\begin{frame}[fragile]{Working with Headers: \texttt{csv.DictReader}}
    \begin{lstlisting}
with open("students.csv", "r", newline="") as f:
    reader = csv.DictReader(f)   # uses first row as keys

    data = []
    for row in reader:
        data.append(row)         # each row is a dictionary

# Each row is now a dict: {'Name': ..., 'Grade': ..., 'Score': ...}
for row in data:
    name  = row["Name"]          # no need to remember column 0
    grade = row["Grade"]         # no need to remember column 1
    score = float(row["Score"])  # convert string to number
    print(f"{name}: {grade} ({score})")
    \end{lstlisting}

    \pause

    \begin{columns}
        \begin{column}{0.48\textwidth}
            \textbf{With \texttt{csv.reader}:}
            \begin{itemize}
                \item \texttt{row[2]} \textrightarrow{} Score
                \item Fragile if columns move
            \end{itemize}
        \end{column}
        \begin{column}{0.48\textwidth}
            \textbf{With \texttt{csv.DictReader}:}
            \begin{itemize}
                \item \texttt{row["Score"]} \textrightarrow{} Score
                \item \textcolor{marigold}{\textbf{Robust and readable!}}
            \end{itemize}
        \end{column}
    \end{columns}
\end{frame}

% -----------------------------------------------
\begin{frame}[fragile]{Try It: Read Your First CSV}
    \begin{center}
        \Large \textbf{Activity Time!}
    \end{center}

    \vspace{1em}

    \textbf{Create a small CSV about your favourite movies:}
    \begin{enumerate}
        \item Open a text editor and create \texttt{movies.csv} with the content:
        \begin{lstlisting}[language={},frame=single,numbers=none]
Title,Year,Rating
Interstellar,2014,9.5
The Matrix,1999,8.7
Coco,2017,8.4
        \end{lstlisting}
        \item Write a Python script that reads it with \texttt{csv.DictReader}
        \item Print each movie title and its rating
        \item Add a 4th movie of your choice and re-run the script
    \end{enumerate}

    \vspace{0.5em}

    \textbf{Bonus:} Print only movies released after the year 2000.
\end{frame}

% -----------------------------------------------
\begin{frame}{Computing Statistics: Min and Max}
    \textbf{Strategy: scan every row, keep track of the best value seen so far}

    \vspace{1em}

    \onslide<2->{
    \textbf{Finding the maximum (pseudocode):}
    \begin{itemize}
        \item<2-> Start with \texttt{current\_max = very small number} (or the first value)
        \item<3-> Loop through every row in the dataset
        \item<4-> Read the value in the column you care about
        \item<5-> Convert it to a number with \texttt{float()}
        \item<6-> If this value \textbf{is greater than} \texttt{current\_max}, update \texttt{current\_max}
        \item<7-> After the loop, \texttt{current\_max} holds the answer
    \end{itemize}
    }

    \vspace{1em}

    \onslide<8->{
    \textbf{Finding the minimum:} same idea - start large, update when smaller.

    \vspace{0.5em}

    \textcolor{marigold}{\textbf{Key insight:}} This is exactly the same loop pattern you used to find the darkest or brightest pixel in an image!
    }
\end{frame}

% -----------------------------------------------
\begin{frame}[fragile]{Computing Max: Code}
    \begin{lstlisting}
import csv

with open("students.csv", "r", newline="") as f:
    reader = csv.DictReader(f)
    data = list(reader)

# ---- find max score ----
current_max = float(data[0]["Score"])   # start with first value
max_student = data[0]["Name"]

for row in data:
    score = float(row["Score"])
    if score > current_max:
        current_max = score
        max_student = row["Name"]

print("Highest score:", current_max, "by", max_student)

    \end{lstlisting}
\end{frame}

% -----------------------------------------------
\begin{frame}[fragile]{Computing the Average}
    \begin{lstlisting}
with open("students.csv", "r", newline="") as f:
    reader = csv.DictReader(f)
    data = list(reader)

total = 0.0
count = 0

for row in data:
    total = total + float(row["Score"])
    count = count + 1


if count > 0:              # Guard against empty files
    average = total / count
    print(f"Average score: {average:.2f}")
else:
    print("No data found.")
    \end{lstlisting}
\end{frame}


% -----------------------------------------------
\begin{frame}{Filtering Rows}
    \begin{center}
        \textbf{Like image filters, we only \textcolor{stfxblue}{keep rows that pass a test}}
    \end{center}

    \vspace{0.5em}

    \begin{center}
        \begin{tikzpicture}[node distance=0.6cm and 1.4cm,
                            box/.style={draw, rounded corners, minimum width=2.2cm,
                                        minimum height=0.55cm, font=\small},
                            arrow/.style={->, thick}]

            % Input column
            \node[box, fill=blue!15]  (r0) {Alice, A, 92};
            \node[box, fill=blue!15, below=of r0] (r1) {Bob, B, 78};
            \node[box, fill=blue!15, below=of r1] (r2) {Carol, A, 95};
            \node[box, fill=blue!15, below=of r2] (r3) {David, C, 64};
            \node[box, fill=blue!15, below=of r3] (r4) {Eve, B, 81};
            \node[above=0.2cm of r0, font=\bfseries\small] {All rows};

            % Filter box
            \node[box, align=center, fill=marigold!40, right=1.6cm of r2, minimum width=2.8cm,
                  minimum height=1.2cm] (filt) {\textbf{Condition:}\\\texttt{Score >= 80}};

            % Output column
            \node[box, fill=green!20, right=1.6cm of filt, yshift=0.6cm]  (o0) {Alice, A, 92};
            \node[box, fill=green!20, below=of o0] (o1) {Carol, A, 95};
            \node[box, fill=green!20, below=of o1] (o2) {Eve, B, 81};
            \node[above=0.2cm of o0, font=\bfseries\small] {Filtered rows};

            % Arrows in
            \draw[arrow] (r0.east) -- (filt.west);
            \draw[arrow] (r1.east) -- (filt.west);
            \draw[arrow] (r2.east) -- (filt.west);
            \draw[arrow] (r3.east) -- (filt.west);
            \draw[arrow] (r4.east) -- (filt.west);

            % Arrows out
            \draw[arrow, green!60!black] (filt.east) -- (o0.west);
            \draw[arrow, green!60!black] (filt.east) -- (o1.west);
            \draw[arrow, green!60!black] (filt.east) -- (o2.west);
        \end{tikzpicture}
    \end{center}

    \vspace{0.3em}

    \onslide<2->{
    \textbf{The pattern:} build a new list containing only the rows we want.
    }
\end{frame}

% -----------------------------------------------
\begin{frame}[fragile]{Filtering: Keep Rows Above a Threshold}
    \textbf{Example: keep only students with score $\geq$ 80}

    \begin{lstlisting}
with open("students.csv", "r", newline="") as f:
    data = list(csv.DictReader(f))
passing = []
for row in data:
    score = float(row["Score"])   # convert string -> number first!
    if score >= 80:
        passing.append(row)

for row in passing:
    print(f"  {row['Name']:10s}  {row['Score']}")
    \end{lstlisting}

    \pause
    \textbf{Output:}
    \begin{lstlisting}
Students with score >= 80 (3 found):
  Alice       92
  Carol       95
  Eve         81
    \end{lstlisting}
\end{frame}

% -----------------------------------------------
\begin{frame}[fragile]{Filtering: Keep Rows Matching a Value}
    \textbf{Example: keep only students with grade \texttt{"A"}}

    \begin{lstlisting}
with open("students.csv", "r", newline="") as f:
    data = list(csv.DictReader(f))
a_students = []
for row in data:
    if row["Grade"] == "A":       # string comparison, no conversion needed
        a_students.append(row)
print(f"A-grade students ({len(a_students)} found):")
for row in a_students:
    print(f"  {row['Name']}")
    \end{lstlisting}

    \pause

    \textbf{Output:}
    \begin{lstlisting}
A-grade students (2 found):
  Alice
  Carol
    \end{lstlisting}
\end{frame}

% -----------------------------------------------
\begin{frame}[fragile]{Sorting Data}
    \textbf{Use Python's built-in \texttt{sorted()} with a \textcolor{stfxblue}{\texttt{key=}} argument:}

    \begin{lstlisting}
with open("students.csv", "r", newline="") as f:
    data = list(csv.DictReader(f))

# Sort by Score, highest first
by_score = sorted(data, key=lambda row: float(row["Score"]),reverse=True)

print("Ranked by score:")
for i, row in enumerate(by_score):
    print(f"  {i+1}. {row['Name']:10s} {row['Score']}")

# Sort alphabetically by Name (ascending)
by_name = sorted(data, key=lambda row: row["Name"])

print("\nAlphabetical order:")
for row in by_name:
    print(f"  {row['Name']}")
    \end{lstlisting}

    \pause

    \begin{itemize}
        \item \texttt{lambda row: ...} is a small, one-line function used as the sort key
        \item \texttt{reverse=True} flips the order from ascending to descending
    \end{itemize}
\end{frame}

% -----------------------------------------------
\begin{frame}[fragile]{Try It: Class Grade Analyzer}
    \begin{center}
        \Large \textbf{Activity Time!}
    \end{center}

    \vspace{0.5em}

    \textbf{Create \texttt{grades.csv} and analyse it:}
    \begin{lstlisting}[language={},frame=single,numbers=none]
Name,Score
Alice,88
Bob,74
Carol,95
David,61
Eve,82
Frank,90
Grace,73
    \end{lstlisting}

    \textbf{Write a script that prints:}
    \begin{enumerate}
        \item The \textcolor{stfxblue}{highest} and \textcolor{stfxblue}{lowest} score (with the student's name)
        \item The \textcolor{marigold}{class average}
        \item A list of students whose score is \textcolor{stfxblue}{above the average}
        \item The class list sorted from \textcolor{marigold}{highest to lowest} score
    \end{enumerate}
\end{frame}

% -----------------------------------------------
\begin{frame}[fragile]{Writing a CSV File}
    \begin{lstlisting}
results = [
    ["Name", "Score", "Grade"],    # header row
    ["Alice", 92, "A"],
    ["Bob",   78, "B"],
    ["Carol", 95, "A"],
    ["David", 64, "C"],
    ["Eve",   81, "B"],
]

with open("output.csv", "w", newline="") as f:
    writer = csv.writer(f)
    for row in results:
        writer.writerow(row)

# You can also write all rows at once:
with open("output2.csv", "w", newline="") as f:
    writer = csv.writer(f)
    writer.writerows(results)   # note: writerows (plural)
    \end{lstlisting}

    \pause

    \texttt{writerow} takes one list; \texttt{writerows} takes a list of lists

\end{frame}

% -----------------------------------------------
\begin{frame}[fragile]{A Complete Data Pipeline}

    \begin{lstlisting}
# 1. LOAD
with open("students.csv", "r", newline="") as f:
    data = list(csv.DictReader(f))
# 2. PROCESS -- compute average score
total = sum(float(row["Score"]) for row in data)
average = total / len(data)
# 3. FILTER -- keep only above-average students
above_avg = [row for row in data if float(row["Score"]) > average]
# Sort filtered results highest first
above_avg.sort(key=lambda row: float(row["Score"]), reverse=True)
# 4. SAVE
with open("above_average.csv", "w", newline="") as f:
    writer = csv.DictWriter(f, fieldnames=["Name", "Grade", "Score"])
    writer.writeheader()
    writer.writerows(above_avg)

print(f"Average: {average:.2f}")
print(f"Saved {len(above_avg)} above-average students.")
    \end{lstlisting}
\end{frame}

% -----------------------------------------------
\begin{frame}{Quick Self-Check}
    \textbf{Can you answer these?}

    \vspace{1em}

    \begin{enumerate}
        \item What does CSV stand for, and what makes it different from a regular text file?
        \item What is the difference between \texttt{csv.reader} and \texttt{csv.DictReader}?
        \item If \texttt{row = ['Alice', 'A', '92']}, what is \texttt{row[2]}?
        \item Why do we need \texttt{float(row["Score"])} before comparing scores?
        \item How does filtering rows from a CSV differ from filtering pixels in an image?
    \end{enumerate}
\end{frame}

% -----------------------------------------------
\begin{frame}{Real-World Applications}
    \textbf{What you have learned is used in:}

    \vspace{1em}

    \begin{itemize}
        \item \textbf{Data science:} Pandas, NumPy, and Jupyter notebooks all start from CSV
        \item \textbf{Finance:} Stock prices, transaction logs, expense reports
        \item \textbf{Sports analytics:} Player statistics, game-by-game results
        \item \textbf{Scientific research:} Experimental measurements, survey results
        \item \textbf{Web scraping:} Save scraped data to CSV for later analysis
        \item \textbf{Machine learning:} Training datasets are almost always tabular
        \item \textbf{Business intelligence:} Sales reports, customer data, A/B test results
    \end{itemize}

\end{frame}

\end{document}
